\chapter{PENDAHULUAN}
\label{chap:pendahuluan}

% Ubah bagian-bagian berikut dengan isi dari pendahuluan

\section{Latar Belakang}
\label{sec:latarbelakang}

Pada dasarnya, sistem tiket elektronik tidak hanya meningkatkan kemudahan dalam
pembelian dan penjualan tiket, tetapi juga membantu dalam manajemen tiket serta
mengurangi risiko kehilangan tiket \parencite{qteishat2014}. Meskipun demikian,
terdapat beberapa permasalahan dalam sistem tiket saat ini. Pertama, sistem ini
bersifat terpusat sehingga harus cukup kuat dalam membangun kepercayaan
pengguna. Dalam sistem terpusat, biasanya terdapat pengendali yang memproses
dan menyimpan data serta menyediakan layanan, yang umumnya berasal dari pihak
ketiga tepercaya \parencite{nakamoto2008bitcoin}. Pihak ketiga ini memang
berperan penting dalam menjaga integritas dan keamanan data, namun juga dapat
menjadi titik lemah apabila tidak andal atau menjadi target serangan siber.
Ketergantungan pada penyedia layanan tertentu juga dapat menyebabkan gangguan
layanan ketika terjadi kesalahan atau serangan, sehingga menciptakan kerentanan
dalam transaksi yang seharusnya aman.

Salah satu solusi yang ditawarkan untuk mengatasi permasalahan ini adalah
adopsi teknologi blockchain, yang menawarkan pencatatan transaksi secara
terdesentralisasi dan transparan \parencite{aventus2018}. Teknologi blockchain
memungkinkan setiap transaksi dicatat dalam blok yang tidak dapat dimodifikasi
dan saling terhubung, memberikan lapisan keamanan tambahan serta memungkinkan
pengguna memverifikasi keabsahan tiket secara mandiri. Dengan demikian,
blockchain menjadi solusi potensial untuk mengatasi isu transparansi, keamanan,
dan integritas data dalam sistem tiket elektronik. Namun, biaya transaksi dan
skalabilitas masih menjadi tantangan dalam implementasi blockchain, terutama
pada jaringan utama Ethereum. Oleh karena itu, dibutuhkan pendekatan lanjutan
guna mengeksplorasi solusi yang lebih efisien dan efektif.

Lebih lanjut, blockchain juga mengubah paradigma kepemilikan aset digital
melalui pengenalan Non-Fungible Token (NFT), yang membuka peluang baru dalam
transaksi dan koleksi digital. Mengingat keterbatasan jaringan utama Ethereum,
solusi Layer-2 seperti Optimism, Polygon zkEVM, dan Arbitrum muncul sebagai
alternatif yang lebih efisien untuk mengimplementasikan NFT. Teknologi Layer-2
memungkinkan transaksi lebih cepat dan murah tanpa mengorbankan keamanan
jaringan utama. Dalam konteks ini, setiap NFT tetap mempertahankan
karakteristik uniknya seperti ID global, kemampuan transfer, serta metadata,
ditambah dengan keunggulan berupa biaya gas rendah dan throughput tinggi.

Standar NFT pada Layer-2 mendukung interoperabilitas dengan ekosistem Ethereum
yang lebih luas, memungkinkan perpindahan aset antara Layer-2 dan jaringan
utama melalui jembatan lintas-rantai yang aman. Ini membuka peluang baru bagi
kreator dan kolektor untuk berdagang secara aman dan efisien. Transparansi
dalam pelacakan kepemilikan dan riwayat transaksi juga dapat meningkatkan
kepercayaan pengguna. Meskipun tantangan seperti hak kekayaan intelektual dan
kompleksitas interoperabilitas masih ada, adopsi NFT pada Layer-2 terus
berkembang dan menunjukkan potensinya sebagai solusi skalabilitas yang layak.
Penggunaan Layer-2 tidak hanya menjawab masalah biaya dan skalabilitas, tetapi
juga membuka kemungkinan baru yang sebelumnya tidak dapat dijangkau secara
ekonomi. Hal ini menunjukkan bahwa teknologi blockchain berpotensi besar dalam
menciptakan sistem kepemilikan digital yang inklusif, efisien, dan
berkelanjutan.

\section{Permasalahan}
\label{sec:permasalahan}

Permasalahan penelitian ini dirumuskan sebagai berikut:
\begin{enumerate}
      \item Bagaimana solusi Layer-2 berbasis \textit{zero-knowledge rollup} seperti
            \textit{Polygon zkEVM} dapat mengatasi keterbatasan skalabilitas dan biaya pada
            sistem tiket konser berbasis \textit{blockchain} Layer-1 Ethereum?
      \item Bagaimana integrasi sistem \textit{Layer-2} dengan \emph{smart contract} NFT
            dapat meningkatkan efisiensi transaksi, menurunkan biaya gas, serta tetap
            menjaga keamanan dan keandalan dalam pengelolaan tiket konser digital?
\end{enumerate}

\section{Tujuan}
Tujuan penelitian ini adalah sebagai berikut:
\begin{enumerate}
      \item Menganalisis bagaimana solusi Layer-2 berbasis \textit{zero-knowledge rollup},
            seperti \textit{Polygon zkEVM}, dapat meningkatkan skalabilitas sistem tiket
            konser berbasis \textit{blockchain} Layer-1 Ethereum, terutama melalui
            pengurangan waktu konfirmasi dan biaya transaksi.
      \item Mengembangkan solusi integrasi sistem Layer-2 yang efisien untuk pengelolaan
            tiket konser digital berbasis \textit{smart contract NFT}, dengan fokus pada
            peningkatan efisiensi, pengurangan biaya transaksi, dan penguatan keamanan
            data.
\end{enumerate}

\section{Batasan Masalah}
\label{sec:batasanmasalah}

Dalam penelitian ini, terdapat beberapa batasan masalah atau ruang lingkup
sebagai berikut:

\begin{enumerate}
      \item Penelitian dibatasi pada praktik penjualan tiket dalam industri konser, tidak
            mencakup acara lain seperti olahraga atau festival seni. Fokus ini untuk
            memberikan analisis mendalam terkait masalah scalping dan penerapan teknologi
            \textit{blockchain} pada konteks konser.

      \item Penelitian hanya menggunakan teknologi \textit{blockchain} dan \textit{NFT}
            sebagai solusi utama untuk permasalahan tiket konser. Teknologi lain seperti
            sistem pembayaran alternatif atau teknologi non-blockchain tidak dibahas dalam
            penelitian ini.

      \item Penelitian menitikberatkan pada aspek keamanan data dan transparansi dalam
            penerapan teknologi \textit{blockchain} dan \textit{NFT} pada tiket konser.
            Aspek sosial dan ekonomi dari penggunaan teknologi ini tidak menjadi fokus
            utama.

      \item Penelitian difokuskan pada penggunaan jaringan \textit{Polygon zkEVM} sebagai
            solusi \textit{Layer-2} dalam pengembangan platform tiket konser berbasis
            \textit{blockchain} dan \textit{NFT}. Solusi \textit{Layer-2} lain, seperti
            \textit{Optimism} atau \textit{Arbitrum}, tidak dianalisis atau dibandingkan.
\end{enumerate}
\section{Manfaat Penelitian}
\label{sec:manfaatpenelitian}

Beberapa manfaat yang diharapkan dari penelitian ini meliputi:

\begin{enumerate}
      \item Implementasi teknologi \textit{blockchain} dapat mengurangi risiko pemalsuan
            tiket dan praktik scalping, sehingga memberikan perlindungan yang lebih baik
            bagi konsumen.

      \item Sistem terdesentralisasi meningkatkan transparansi dalam proses verifikasi
            keaslian tiket, sehingga pengguna dapat lebih yakin terhadap keaslian tiket
            yang mereka beli.

      \item Dengan menurunkan biaya transaksi dan meningkatkan keandalan dalam proses
            pembelian tiket, pengalaman pengguna dapat ditingkatkan, sehingga lebih banyak
            orang dapat menikmati acara tanpa khawatir akan harga yang tidak wajar.
\end{enumerate}
